\section{Experiments}
\label{sec:experiments}

\subsection{Data and protocol}
\label{sec:protocol}
\paragraph{Leaderboards.}
We use two public leaderboards, both as frozen snapshots. \textbf{HELM MMLU} v1.13.0 (release of 2025-01-10): exact-match accuracy of $M=79$ models on the $n=57$ MMLU subjects \citep{hendrycks2020mmlu,liang2022holistic}; the subjects have 100 to 1534 test items. \textbf{Open LLM Leaderboard} (v2, \texttt{contents} dataset, commit \texttt{9c09a7c}, last updated 2025-03-20): $n=6$ tasks (IFEval, BBH, MATH Lvl~5, GPQA, MuSR, MMLU-Pro), where BBH \citep{suzgun2022bbh} and the others are reported as the leaderboard's own normalised scores (rescaled against the random baseline). Of 4576 rows we drop 79 duplicate model names (keeping the best average), leaving 4497 models, and exclude those the leaderboard flags, leaving 4496 models and 4495 targets (the model that is first under the mean is not a target). Checksums, licences and download commands are in the reproduction notes.

\paragraph{Targets, rules, caps.}
Every model is taken as the target in turn, the others being fixed, so the unit of analysis is the (target, rule, cap) triple. Rules and caps are those of Section~\ref{sec:method}. Because the exact mixed-integer programmes for the rank rules are expensive on 4496 models, the OLLB rank-rule grids (Borda, plurality, Condorcet winner) use a rank-stratified sample of 57 targets (the 20 best-ranked models and 40 at log-uniformly spaced ranks); the score rules use all 4495 targets. Every rule, including Copeland and the trimmed mean, is run on a top-500 slice (the 500 best models by mean, with a rank-stratified sample of 116 targets; Copeland on 30 targets at cap 100 and on all 116 at base $+10$). Medians over a sample of targets are not comparable with medians over all targets: cross-rule statements on OLLB use the same targets (paired ratios, the top-500 slice).

\paragraph{Summaries and uncertainty.}
Costs are medians over targets, always with the share of infeasible targets alongside (infeasible targets are never given a cost); confidence intervals are percentile bootstrap intervals over targets (seed 0, 1000 resamples). The task-level and evaluation-noise uncertainty of a single target's cost is studied separately (Section~\ref{sec:robust}).


\subsection{How much investment, by aggregation rule}
\label{sec:byrule}
Figure~\ref{fig:byrule} and Appendix Table~\ref{tab:headline} give the minimum investment to first place on HELM MMLU with each of the 78 models that are not already first under the mean rule as target, with the cap at 100 (every task can be raised to the ceiling) and with the cap at the best existing score on each subject. With the cap at 100 every target can reach first place under every rule, so all medians below are over all 78 targets.

\paragraph{The rule changes the price by an order of magnitude.}
Under the mean the median target must raise 16 of the 57 subjects (95\% CI over targets $[13,18.5]$) by $L_1=696$ points in total ($[556,875]$). Under plurality the median is $L_1=47$ points on 9 subjects; under the minimum, 98 points on 11.5; under the median, 112 points on 21; for becoming Condorcet winner, 213 points on 29; under Borda, 534 points on 47 (Appendix Table~\ref{tab:headline}). Paired by target, plurality, the minimum and the median cost 0.07, 0.15 and 0.16 of the mean.

\paragraph{Points and tasks do not rank the rules the same way.}
Borda costs 0.78 of the mean in points but needs the most subjects (47, 82\% of the 57): the Borda credit a target gains on one subject is bounded by the number of models it passes. The median and the minimum need few points but touch 21 and 11.5 subjects, more than plurality (9). The geometric, harmonic and power means lower the points (paired ratios 0.98 to 0.77 from $p=0.5$ to $p=-5$) but not the number of subjects (15.5--16; Appendix Figure~\ref{fig:power}). With the cap at the best existing score, the order of the rules in $L_1$ is almost unchanged (Kendall's $\tau$ between caps, median over pairs of HELM caps: $0.85$, range $0.63$--$0.93$), the order in $K$ is not (median $\tau=0.63$, range $0.19$--$0.87$); in $K$ the mean goes from 16 to 37 subjects and plurality from 9 to 18.

\paragraph{Borda and the Condorcet requirement are structural, not accidental.}
For the median target the Condorcet-winner plan touches 29 subjects at every cap, the majority threshold $\lfloor57/2\rfloor+1$; Copeland, bounded above by it, has the same $K$ and almost the same $L_1$ (213 versus 213 points at 100; 254 versus 259 at the best score). Five weighted means with Dirichlet$(1,\ldots,1)$ weights are cheaper than the plain mean: median $K$ 7--9 (mean: 16), paired $L_1$ ratio 0.36--0.42 (Appendix Figure~\ref{fig:weighted}), more than the choice among geometric, harmonic or power means with $p\ge-2$.

\paragraph{One target, five plans.}
For the model ranked 5 under the mean (Claude 3 Opus), with each subject capped at $+10$ points, an $L_1$-optimal plan costs 154 points on 16 subjects under the mean, 6.2 points on 5 subjects under the median, 5.0 points on 10 under plurality, 43.7 on 46 under Borda and 24.5 on 27 for the Condorcet winner (Appendix Figure~\ref{fig:example_plan}; optimal plans are not unique, and the $L_1$-optimal plan of the ordinal rules touches more subjects than their $K$-minimal one). The plans do not transfer: the mean plan makes the model first under the median and plurality but fourth under Borda; the median and plurality plans leave it fifth under the mean.

\begin{figure}[t]
\centering
\includegraphics[width=\textwidth]{fig_cost_by_rule_helm.pdf}
\caption{Minimum targeted investment to first place by aggregation rule and cap on each task's score (HELM MMLU, 78 targets, 57 subjects). Top: subjects touched $K$ (minimised); bottom: total points $L_1$ (minimised separately). Dots are targets; bar: interquartile range; circle: median, over the targets for which first place is reachable (feasible/79 printed at the right; infeasible targets are not drawn and not counted as zero). $^*$Condorcet winner: an upper bound for Condorcet-consistent rules. Caps: 100; best existing score on each subject; target's score $+20$.}
\label{fig:byrule}
\end{figure}


\subsection{Cost against the target's current standing}
\label{sec:rank}
The cost grows with the distance to first place (Figure~\ref{fig:rank}). Across targets, $L_1$ is almost a monotone function of the target's rank under the rule itself (Spearman $\rho=0.99$ for the median, $0.88$ for the minimum, $0.61$ for plurality, cap 100) and of its score shortfall to the rule's leader (median $\rho=0.998$ over 14 rules; lowest $0.88$). The slope of $\log L_1$ on $\log(\text{rank})$ at cap 100 is $1.13$ ($[0.92,1.34]$) for the mean, $1.91$ ($[1.71,2.07]$) for the median, $1.59$ for the minimum and $3.29$ ($[2.21,4.11]$) for plurality: the rules that are cheapest near the top become expensive faster with rank (the ranks span only 2--79, so we report these as descriptive slopes). Targets close to the top are cheap: for the mean, the rank-2 model needs 4.0 points on 1 subject, ranks 3--5 a median of 42 points on 1 subject, ranks 6--10 177 points on 5 subjects, ranks 11--30 435 points on 10. The headline medians of Appendix Table~\ref{tab:headline} are dominated by the 49 targets ranked 31 or below, which no developer would pick; we therefore always show cost against rank.

\begin{figure}[t]
\centering
\includegraphics[width=\textwidth]{fig_cost_vs_rank_helm.pdf}
\caption{Investment against the target's rank under the mean (78 targets), cap 100 (left) and best existing score (right). Per rank bin: median $K$ (top) and $L_1$ (bottom, symmetric log) with 95\% bootstrap interval of the median (1000 resamples, seed 0). Bins hold 3--32 targets; the interval of the lowest bins is wide. $^*$Condorcet winner.}
\label{fig:rank}
\end{figure}

\subsection{Caps and feasibility: when no plan exists}
\label{sec:feas}
Caps decide whether the answer is a cost or ``impossible''. With the three permissive caps (100; best existing score $+2$; $+0$) every target is feasible under every rule. With a cap tied to the target's own score the picture changes (Figure~\ref{fig:feas}): the share of targets for which first place is unreachable at $+10$ points is 0.63 (mean), 0.49 (median), 0.27 (plurality), 0.51 (Copeland), 0.68 (Borda) and 0.82 (minimum), at $+5$ points 0.86, 0.79, 0.63, 0.79, 0.86 and 0.91, and at $+20$ points 0.28, 0.15, 0.05, 0.17, 0.37 and 0.56. With a $+10$ gain under the mean, no model ranked 31 or lower can reach first place, whereas 20 of the models ranked 11--30 can, with a median of 44.5 subjects. The minimum is the rule that is cheapest when feasible and the one most often infeasible (82\% at $+10$): it demands the target's \emph{worst} subject to exceed the best rival's worst.

\begin{figure}[t]
\centering
\includegraphics[width=\textwidth]{fig_feasibility_helm.pdf}
\caption{(a) Number of the 79 HELM models for which first place is reachable at all, by rule and cap. (b) Share of targets that can reach first place with the cap at the target's score $+20$, by the target's rank under the mean (bins of ten ranks).}
\label{fig:feas}
\end{figure}

\paragraph{Selection effect.}
Medians under a tight cap are over different feasible subsets (e.g.\ at $+10$, 29 targets for the mean and 57 for plurality): such a median is not the cost of the same target set, and a rule with fewer feasible targets looks cheaper because only its easy targets are feasible (the minimum at $+10$: 3 subjects, 9 points, for 14 of 78 targets). For cross-rule claims we therefore use the paired ratio on targets feasible under both rules and the ranking at the permissive caps.

\subsection{Number of sub-datasets}
\label{sec:kn}
Does a majority-type rule force about $n/2$ subjects? We drew random subsets of $n\in\{3,6,12,24,48\}$ HELM subjects (60 draws each) and followed targets at rank 2 and 5 under the mean (Appendix Figure~\ref{fig:kn}). The worst case of the median rule is $K=\lfloor n/2\rfloor+1$ (Appendix~\ref{app:bounds}) and of the minimum $K=n$, but targets near the top come nowhere near it: at $n=48$ the rank-5 target needs $K=3$ under the median (bound: 25), 4 under plurality and 1 under the minimum. With the cap at $+10$ the mean and Borda grow in proportion to $n$ ($K=1,2,4,7,14$ for the mean and $2,3,4.5,7.5,17$ for Borda at $n=3,\ldots,48$, about $0.3n$). The bound is thus tight for a worst-case target and false for contenders: the mean and Borda need a roughly constant share of the subjects, the median, the minimum and plurality a few.

\subsection{Robustness to evaluation noise and to the task set}
\label{sec:robust}
Scores are measured with error on a finite item set. For the ten models ranked 2--11 under the mean, with the cap at $+10$, we redrew every score from its binomial noise (the item count of each subject; 100 draws, seed 0) and, separately, resampled the 57 subjects with replacement (100 draws), and recomputed the minimum investment under seven rules (Appendix Figure~\ref{fig:robust}). The width of the 5--95\% interval of $L_1$ relative to its median is 0.51 (noise) and 0.76 (task resampling) for the mean, but 2.33 and 1.90 for the median, 1.90 and 2.56 for the minimum, 1.49 and 1.51 for plurality, 1.13 and 1.10 for Borda and 1.08 and 1.01 for the Condorcet winner (median over the ten targets, feasible draws). For the minimum, first place is unreachable in 44\% of the noise draws and 30\% of the task resamples (average over the ten targets). Hence the rules that are cheapest on the observed data are also the least stable: for those rules the 5--95\% interval is wider than the median itself. The margin $\varepsilon=10^{-3}$ is far below the evaluation noise (3.5 points for a subject with the median 173 items at 70\% accuracy), so ``first'' is a statement about the observed table, not about the underlying models.


\subsection{Second leaderboard: the Open LLM Leaderboard}
\label{sec:ollb}
The Open LLM Leaderboard has only $n=6$ tasks and 4495 targets, most of them far from first place (4396 rank below 100 under the mean), so $K$ is small: at cap 100 the median $K$ is at most 3 of 6 tasks for every rule (4 for the Condorcet winner on the top-500 slice); under tight caps the few feasible targets need up to 5 or 6 (Appendix Tables~\ref{tab:headline_ollb} and~\ref{tab:headline_ollb_top500}, Figures~\ref{fig:byrule_ollb} and~\ref{fig:rank_ollb}). Under the mean with the cap at 100, the median target needs 181 points ($[178,182]$) on 2 tasks; the model ranked 2 needs 4.8 points on one task, the three models ranked 3--5 a median of 7.9 points, those ranked 6--10 27.7 points, all on one task (log--log slope of $L_1$ on rank $0.44$, $[0.43,0.44]$, against $1.13$ on HELM).

\paragraph{Only the minimum transfers.}
Paired by target at cap 100 (Table~\ref{tab:transfer}), the minimum costs 0.15 of the mean on HELM, 0.23 on the top-500 slice (116 targets, every rule on the same targets) and 0.25 on all OLLB targets. The saving of the other rules does not carry over: the median rule costs 0.16, 0.34 and 0.60, plurality 0.068, 0.38 and 0.50 (the last on the 57 sampled targets), and Borda 0.78, 0.39 and 0.37; on the top-500 slice Copeland costs 0.27 (27 targets with a finished solve) and the trimmed mean 0.82 (Appendix Table~\ref{tab:headline_ollb_top500}). Borda, which touches 82\% of the subjects on HELM, needs 2--3 of the 6 tasks on OLLB. Kendall's $\tau$ between the rule orders of HELM and OLLB is 0.63 in $L_1$ (17 rules) and 0.46 in $K$ against all OLLB targets, and 0.46 and 0.49 against the top-500 slice; among OLLB settings (caps and variants) the median $\tau$ is 0.68 in $L_1$ and 0.58 in $K$. The order of the rules is thus a property of the leaderboard as well as of the rules.

\begin{table}[t]
\centering
\caption{Does the cost order of the rules carry over from HELM to the Open LLM Leaderboard (OLLB)? Cap 100 (every target feasible). Per target set: median paired ratio $L_1/L_1^{\text{mean}}$ with 95\% bootstrap interval over targets, and median number of tasks touched $K$. $^\dagger$Run on 57 sampled targets only (ratio to the mean on the same 57). $^\ddagger$27 targets with a finished solve out of 30 run. $^*$Condorcet winner: upper bound for Condorcet-consistent rules. Last row: Kendall's $\tau$ between the rule orders ($L_1$ / $K$) of HELM and of the OLLB set, over the rules run in both (including power, weighted and trimmed means not shown).}
\label{tab:transfer}
\scriptsize
\resizebox{\textwidth}{!}{\input{../tables/tab_rule_transfer.tex}}
\end{table}

\paragraph{Feasibility.}
With a gain of at most 10 points per task, first place is unreachable for 97\% of the OLLB targets under the mean (113 of 4495 feasible), 95\% under the median and 90\% under the minimum; with the 95th percentile of the observed gains of fine-tuned models over their base models as the cap (\texttt{emp95}) it is unreachable for 93\% (mean), 97\% (median) and 99\% (minimum). On the top-500 slice, where the targets are closer to the top, the corresponding shares at $+10$ points are 34\% (mean), 31\% (median), 25\% (minimum), 10\% (Borda), 29\% (Copeland), 59\% (Condorcet winner) and 65\% (plurality). The empirical cap is a sensitivity scenario and is not calibrated to HELM (Appendix Figure~\ref{fig:feas_ollb}).

\paragraph{Raw against normalised scores.}
\label{sec:rawnorm}
The Leaderboard reports scores normalised against the random baseline of each task. We re-ran the grids on the raw scores and compared matched targets (cap 100). The cost in raw points over the cost in normalised points has median 0.85 for the mean, 0.86 for the median, 0.91 for the geometric mean and 0.99 for the harmonic mean, but 1.38 for the power mean with $p=-5$ and 1.61 for the minimum; the number of tasks touched is the same for 94\% (mean) and 29\% (minimum) of the targets, and the target's rank under the mean is the same for 1\% of the targets (Appendix Figure~\ref{fig:rawnorm}). The two scales are different units, so the cost ratios are not a measure of bias, but their spread (1.9 times between the mean and the minimum) shows that normalisation changes the price of first place by a factor that depends on the rule and moves the order of the rules. Excluding merged models (3788 targets) changes the paired ratios by at most 0.08 (minimum 0.24, median 0.60, geometric mean 0.85, Borda 0.38, Condorcet 0.59, plurality 0.42).

\paragraph{Evaluation noise.}
For the ten models ranked 2--11 on the top-500 slice (cap base $+10$, 50 draws with binomial noise of each task's item count; seed 0), the width of the 5--95\% interval of $L_1$ relative to its median is 0.52 for the mean, 0.86 for the geometric mean and 0.65 for the Condorcet winner, but 1.83 for the median and 1.59 for Borda (median over targets). First place is unreachable in 52\% of the draws under the minimum and in 30\% under plurality, and the target's rank under the minimum rule is about 248 in the median draw: that rule is very sensitive to the single weakest cell (Appendix Figure~\ref{fig:robust_ollb}).
